\documentclass[11pt]{article}
\usepackage[margin=1in]{geometry}
\usepackage{amsmath,amssymb,amsthm,mathtools}
\usepackage{graphicx}
\usepackage{microtype}
\usepackage[hidelinks]{hyperref}

\newtheorem{theorem}{Theorem}
\newtheorem{lemma}{Lemma}
\newcommand{\Mideal}{\mathcal L^{(1,\infty)}}
\newcommand{\Trw}{\operatorname{Tr}_{\omega}}

\title{A Connes--Dixmier-measurable operator\\that is not Dixmier measurable}
\author{Candidate counterexample for arXiv:math/0608375}
\date{11 August 2026}

\begin{document}
\maketitle

\begin{center}
\fbox{\parbox{0.92\linewidth}{\textbf{Status.} Candidate full negative answer
to Question (i) of Carey--Sukochev, subject to expert review.}}
\end{center}

\section{Source question and result}

Carey and Sukochev ask, immediately after Corollary 6.8 on page 22 of
their survey \cite{CS}, whether the spaces of Dixmier-measurable and
Connes--Dixmier-measurable operators coincide.  Their corollary proves
coincidence on the positive cone of the Dixmier--Macaev ideal
$\Mideal$.

\begin{figure}[ht]
  \centering
  \includegraphics[width=0.94\linewidth]{figures/open_problem_crop.png}
  \caption{The open question in \cite[p.~22]{CS}.}
\end{figure}

We answer the question negatively, already for a diagonal self-adjoint
operator on $\ell_2\oplus\ell_2$.

For a positive compact operator $T$, put
\[
 a_T(t)=\frac{1}{\log(1+t)}\int_0^t\mu_s(T)\,ds,
 \qquad t>0.
\]
If $T=T_+-T_-$ is self-adjoint, its trace profile is
\[
 f_T(t)=a_{T_+}(t)-a_{T_-}(t).
\]
Every Dixmier trace has the form $\Trw(T)=\omega(f_T)$ for a
dilation-invariant generalized limit $\omega$.  A Connes--Dixmier trace uses
$\omega=\gamma\circ M$, where
\[
 (Mf)(t)=\frac1{\log t}\int_1^t f(s)\,\frac{ds}{s},\qquad t>1,
\]
and $\gamma$ is a singular state.  These are precisely the definitions used
in \cite[Sections~5--6]{CS}.

\begin{theorem}\label{thm:main}
There is a compact self-adjoint operator
$B\in\Mideal(\ell_2\oplus\ell_2)$ such that
\begin{enumerate}
\item $\Trw(B)=0$ for every Connes--Dixmier trace;
\item there are Dixmier traces $\operatorname{Tr}_{\omega_0},
\operatorname{Tr}_{\omega_1}$ satisfying
\[
 \operatorname{Tr}_{\omega_0}(B)=0,
 \qquad
 \operatorname{Tr}_{\omega_1}(B)=\frac{2}{\log 2}.
\]
\end{enumerate}
Consequently, the Connes--Dixmier-measurable and Dixmier-measurable
operator spaces in $\Mideal$ do not coincide.
\end{theorem}

\section{Proof intuition}

On logarithmic scale, the operator $M$ is the ordinary Ces\`aro average.
Thus a bounded profile made of longer and longer plateaux can still satisfy
$Mf\to0$ if those plateaux have asymptotic density zero.  Every
Connes--Dixmier trace then sees zero.  On the other hand, an invariant mean
can be obtained by averaging only over the plateaux; it sees their nonzero
height.  A second invariant mean can average over the much longer zero gaps.

The only real constraint is to realize this abstract profile by an operator.
In the dyadic rank block $[2^n,2^{n+1})$, prescribe total singular mass
$p_n$.  The singular values are decreasing exactly when
$p_{n+1}\le 2p_n$.  A doubling ramp therefore creates a macroscopic mass
jump without violating monotonicity.  A matching ramp in the negative Jordan
part cancels the jump exactly after a long plateau.  Superexponentially
separated pulses keep both cumulative singular masses $O(n)$ through dyadic
time $n$, which is exactly the $\Mideal$ bound.

\section{Construction}

For $j\ge2$, define
\[
 N_j=2^{2^j},\qquad L_j=2^j,\qquad W_j=2^{2^{j-1}}=\sqrt{N_j}.
\]
In particular, $2^{L_j}=N_j$ and $N_{j+1}=N_j^2$.  The intervals below are
pairwise disjoint.  Define two sequences of dyadic block masses by
\[
 p_n=\begin{cases}
  2^r,&n=N_j-L_j+r\text{ for some }j\ge2, 1\le r\le L_j,\\
  1,&\text{otherwise},
 \end{cases}
\]
and
\[
 q_n=\begin{cases}
  2^r,&n=N_j+W_j+r\text{ for some }j\ge2, 1\le r\le L_j,\\
  1,&\text{otherwise}.
 \end{cases}
\]
For integer ranks $k$ in the dyadic block
$I_n=\{2^n,\ldots,2^{n+1}-1\}$, set
\[
 u_k=\frac{p_n}{2^n},\qquad v_k=\frac{q_n}{2^n}.
\]
Assign any harmless common value at rank zero.  Let
\[
 U=\operatorname{diag}(u_k),\qquad
 V=\operatorname{diag}(v_k),\qquad
 B=U\oplus(-V).
\]

\begin{lemma}\label{lem:ideal}
The sequences $(u_k)$ and $(v_k)$ are decreasing to zero, and
$U,V\in\Mideal$.  Consequently $B\in\Mideal$ and its Jordan parts are
$B_+=U\oplus0$ and $B_-=0\oplus V$.
\end{lemma}

\begin{proof}
The only upward changes in either block-mass sequence are the steps
$1,2,4,\ldots,2^{L_j}$.  Hence always
$p_{n+1}\le2p_n$ and $q_{n+1}\le2q_n$.  Therefore
\[
 \frac{p_{n+1}}{2^{n+1}}\le\frac{p_n}{2^n},\qquad
 \frac{q_{n+1}}{2^{n+1}}\le\frac{q_n}{2^n}.
\]
The ratios tend to zero because every exceptional mass near logarithmic
time $N_j$ is at most $N_j$, whereas the block length is $2^n$.

The excess mass in one complete ramp is
\begin{equation}\label{eq:A}
 A_j:=\sum_{r=1}^{L_j}(2^r-1)
     =2^{L_j+1}-2-L_j=2N_j-2-L_j<2N_j.
\end{equation}
For any $m$, the sum $\sum_{n\le m}p_n$ is its baseline $m+1$ plus
complete earlier ramp excesses and at most one current partial excess.
If the $j$th ramp has begun, then $m\ge N_j-L_j+1\ge N_j/2$.
Moreover, since $N_{i+1}=N_i^2$,
$\sum_{i<j}N_i\le N_j$ (enlarging the harmless constant for the first
few indices).  Thus
\[
 \sum_{n=0}^m p_n\le C(m+1)
\]
with an absolute $C$.  The same argument applies to $(q_n)$, whose ramp is
even later than $N_j$.  At the end of block $I_m$, cumulative singular mass
is therefore $O(m)$, while $\log(1+2^{m+1})\asymp m+1$.  The estimate inside
a block follows by monotonicity.  This is exactly the $\Mideal$ bound.
\end{proof}

\section{Trace-profile analysis}

Let $F=f_B$.  The common baseline block masses cancel.  During the positive
ramp of pulse $j$, the numerator of $F$ increases from zero to $A_j$; it is
constant through the following $W_j$ baseline blocks; the negative ramp has
the same excesses in the same order and decreases it exactly back to zero.
It follows that, in logarithmic coordinates,
\begin{equation}\label{eq:support}
 \operatorname{supp}\bigl(v\mapsto F(2^v)\bigr)
 \subseteq
 \bigcup_{j\ge2}
 [N_j-L_j+1,\,N_j+W_j+L_j+1].
\end{equation}
The profile is uniformly bounded: on the $j$th interval its numerator is
between $0$ and $A_j<2N_j$, and its logarithmic denominator is comparable
with $N_j$.

\begin{lemma}\label{lem:CD}
$MF(t)\to0$ as $t\to\infty$.
\end{lemma}

\begin{proof}
For $x>0$, the change of variables $s=2^v$ gives
\[
 (MF)(2^x)=\frac1x\int_0^x F(2^v)\,dv.
\]
By \eqref{eq:support}, the logarithmic support accumulated up through a
point of order $N_j$ is at most
\[
 \sum_{i\le j}(W_i+2L_i+2)=O(W_j)=O(\sqrt{N_j})=o(N_j).
\]
The same estimate, with the last completed pulse, applies in every gap.
Since $F$ is uniformly bounded, its average on $[0,x]$ tends to zero.
\end{proof}

Every Connes--Dixmier trace on $B$ equals $\gamma(MF)$ for a singular state
$\gamma$.  Lemma~\ref{lem:CD} proves part (1) of
Theorem~\ref{thm:main}.

For the converse separation, consider the plateau intervals
\[
 J_j=[N_j+1,N_j+W_j+1].
\]
On $J_j$ the numerator of $F(2^v)$ is exactly $A_j$, so uniformly for
$v\in J_j$,
\begin{equation}\label{eq:plateau}
 F(2^v)=\frac{A_j}{\log(1+2^v)}
 \longrightarrow \frac{2}{\log2},
\end{equation}
because $A_j/N_j\to2$ and $W_j/N_j\to0$.  The interval lengths
$|J_j|=W_j$ tend to infinity.

Let $\mathcal U$ be a free ultrafilter and, for $g\in L_\infty(\mathbb R_+)$,
define
\[
 \Lambda_1(g)=\lim_{j\to\mathcal U}
   \frac1{|J_j|}\int_{J_j}g(v)\,dv.
\]
This is a generalized limit.  It is translation invariant: translating a
fixed distance changes an interval average by at most the norm of $g$ times
the relative boundary length, which tends to zero.  Hence the functional
\[
 \omega_1(h):=\Lambda_1\bigl(v\mapsto h(2^v)\bigr)
\]
is a dilation-invariant generalized limit on $L_\infty(\mathbb R_+^*)$.
Equation \eqref{eq:plateau} gives
$\omega_1(F)=2/\log2$.

After the negative ramp, $F$ is identically zero until the next positive
ramp.  In particular, for all sufficiently large $j$ it vanishes on
\[
 K_j=[2N_j,3N_j],
\]
because the $j$th pulse ends before $2N_j$ and the next begins near
$N_{j+1}=N_j^2$.  Repeating the interval-ultralimit construction with
$K_j$ gives a dilation-invariant generalized limit $\omega_0$ satisfying
$\omega_0(F)=0$.  The corresponding Dixmier traces disagree on $B$.
This proves part (2) and Theorem~\ref{thm:main}.

\section{Verification and novelty boundary}

The script \texttt{code/verify\_pulses.py} checks, for increasingly enormous
exact integer parameters, the ramp identity \eqref{eq:A}, exact positive--
negative cancellation, the local inequalities $p_{n+1}\le2p_n$ and
$q_{n+1}\le2q_n$, the prefix-growth estimate, and convergence of the plateau
height to $2/\log2$.

A bounded novelty search used the exact source wording and the phrases
``Connes--Dixmier measurable but not Dixmier measurable'', ``non-positive'',
``self-adjoint'', and ``$\mathcal M_{1,\infty}$''.  Sukochev--Usachev--Zanin
\cite{SUZ} prove that the two trace classes differ on
$\mathcal M_{1,\infty}$, but their positive measurable-operator separation
uses a different Marcinkiewicz ideal; their paper reiterates positive
coincidence on $\mathcal M_{1,\infty}$.  Semenov--Sukochev--Usachev--Zanin
\cite[Theorem~7.7]{SSUZ} later drop positivity on the strictly smaller weak
trace ideal $\mathcal L_{1,\infty}$ and explicitly label that scope.  No source
found in this search states the arbitrary-operator separation above in the
original $\Mideal$.

The result addresses only Question (i) of \cite{CS}; it gives no description
for the maximally invariant or all-symmetric-functional classes in Questions
(ii) and (iii).

\section{Human-review recommendation}

Review as a high-priority candidate full counterexample.  The most useful
checks are: (a) the uniform prefix bound in Lemma~\ref{lem:ideal}; (b) the
support-density estimate in Lemma~\ref{lem:CD}; and (c) that the interval
ultralimits meet exactly the source definition of dilation-invariant states.

\begin{thebibliography}{9}
\bibitem{CS}
A.~L. Carey and F.~A. Sukochev,
\emph{Dixmier traces and some applications to noncommutative geometry},
Russian Math. Surveys 61 (2006), 1039--1099; arXiv:math/0608375.

\bibitem{SUZ}
F.~Sukochev, A.~Usachev, and D.~Zanin,
\emph{On the distinction between the classes of Dixmier and
Connes--Dixmier traces},
Proc. Amer. Math. Soc. 141 (2013), 2169--2179; arXiv:1210.3397.

\bibitem{SSUZ}
E.~Semenov, F.~Sukochev, A.~Usachev, and D.~Zanin,
\emph{Banach limits and traces on $\mathcal L_{1,\infty}$},
Adv. Math. 285 (2015), 568--628; arXiv:1612.04509.
\end{thebibliography}

\end{document}
